\uuid{W55Z}
\titre{Pour tout/il existe}
\auteur{bodin ; hanani ; mzari}
\organisation{exo7}
\licence{CC BY-NC-SA 4.0}
\datecreate{2020-01-20}
\niveau{L1}
\module{Algèbre}
\chapitre{Logique, ensemble, raisonnement}
\sousChapitre{Absurde et contraposée}
\difficulte{1}
\theme{Raisonnement, pour tout, il existe}
\exo7qcmid{lille-1.logique.022}

\contenu{
\begin{question}
\qclassification{100.03,100.04}
\qtags{motcle=pour tout/il existe}
Chercher un contre-exemple à une assertion du type
"$\forall x \in E$ l'assertion $P(x)$ est vraie" revient à prouver l'assertion :
\begin{answers}
    \bad{$\exists! x \in E \quad$ l'assertion $P(x)$ est fausse.}

    \good{$\exists x \in E \quad$ l'assertion $P(x)$ est fausse.}

    \bad{$\forall x \notin E \quad$ l'assertion $P(x)$ est fausse.}

    \bad{$\forall x \in E \quad$ l'assertion $P(x)$ est fausse.}
\end{answers}
\begin{explanations}
Un contre-exemple, c'est trouver un $x$ qui ne vérifie pas $P(x)$. (Rien ne dit qu'il est unique.)
\end{explanations}
\end{question}
}
